\section{Related Work}
\label{sec:related}

\textbf{On-policy distillation for agents.}
ReOPD identifies a prefix trap: histories closer to the student's distribution can make the teacher's targets less reliable, motivating a curriculum over replayed teacher prefixes~\citep{reopd2026}.
Our compared methods address related trajectory-level difficulties: TCOD-F2B schedules rollout depth, Guided-OPD allocates teacher-controlled turns, and FutureBridge validates teacher bridges from high-disagreement states~\citep{tcod2026,guidedopd2026,futurebridge2026}.
These methods retain the original scorer.
\method{} repairs its token scores before student training, a change that can be combined with student-side curricula.

\textbf{Context and self-distillation.}
SDFT and Self-Distilled Reasoner use demonstration-conditioned self-teachers to supervise student rollouts~\citep{sdft2026,opsd2026}.
SDPO instead conditions on textual feedback or successful attempts~\citep{sdpo2026}, while OPCD internalizes experiential knowledge and system prompts through on-policy context distillation~\citep{opcd2026}.
CANON selects privileged solutions by consensus without gold labels~\citep{canon2026}.
These works establish that useful distillation signals need not come from task rewards.
\method{} uses teacher takeovers to repair an external scorer, then freezes it for a separate student's OPD.

\textbf{Teacher quality and repair.}
\citet{li2026bestval} show that a higher-scoring teacher need not transfer better: compatible reasoning patterns and new capabilities matter.
\citet{rethinkingopsd2026} find that privileged context can suppress self-correction and degrade thinking models.
These results motivate checking what the teacher's local signal teaches, beyond its standalone performance.
Distillation Traps and Guards directly calibrates an external teacher with utility and calibration rewards~\citep{trapsguards2026}.
\method{} instead fits takeover-informed token targets without task reward and evaluates the repaired scorer through target fidelity and downstream completion.
\Cref{app:compare,app:related} give the detailed comparisons and extended discussion.
